\chapter{Phase 4 Foundations: Literature Review, Canonical Dictionary, and Epistemic Taxonomy}
\label{ch:phase4_foundations}

% EDITORIAL QUESTION: How do we establish a single, rigorous canonical unit standard that harmonizes conflicting literature representations and eliminates the loose "gauge gap" misconception?
\section{Motivation, Context, and Normalization Obstructions}
\label{sec:phase4_motivation}

The Kitaev honeycomb model under a weak magnetic field $\bm{h} \parallel [111]$ has become the canonical paradigm for non-Abelian topological quantum spin liquids (QSLs)~\cite{Kitaev2006}. However, independent theoretical and numerical explorations across two decades have accumulated disparate operator normalizations, exchange definitions, and heuristic approximations. In particular, two persistent obstructions hinder direct comparison:
\begin{enumerate}
    \item \textbf{Unit and Normalization Ambiguity:} Some works employ spin angular momentum operators $\bm{S} = \frac{1}{2}\bm{\sigma}$ with exchange coupling $J$, others set $K = 1/2$, and others use Pauli operators $\bm{\sigma} \in \{\pm 1\}$ with $K \equiv 1.0$. Because the chiral Majorana gap opens at third order in the magnetic field ($\Delta \propto h^3 / K^2$), a factor of $2$ in the spin definition induces a factor of $2^3 / 2^2 = 2$ in the gap or an overall factor of $4$ in the physical excitation spectrum.
    \item \textbf{The "Gauge Gap" Misconception:} In informal literature discussions, the energy required to invert a $\mathbb{Z}_2$ link variable $u_{ij} \to -u_{ij}$ is frequently called the "gauge excitation gap." This terminology is physically imprecise. In the microscopic spin Hilbert space $\mathcal{H}_{\mathrm{phys}}$, a single link flip is not a physical state; it creates an isolated pair of non-local $\mathbb{Z}_2$ vortices (visons). In the unprojected Hartree-Fock-Bogoliubov (HFB) state space, an isolated link flip does not provide a stationary fixed point of the self-consistent map $\Phi$. When iterated to convergence, the surrounding matter and gauge bonds undergo dynamical screening.
\end{enumerate}

To establish absolute mathematical rigor, Phase 4 strictly enforces a \textbf{single canonical Pauli normalization} across all analytical formulas, code modules, and numerical benchmarks.

% EDITORIAL QUESTION: What is the exact mathematical dictionary mapping alternative literature conventions to the canonical standard?
\section{The Single Canonical Pauli Standard and Universal Dictionary}
\label{sec:phase4_dictionary}
% CLAIM: K-024 (provenance: doc/rnd/phase_4/analytical_prefactor_derivation.md, test: test/unit/test_phase4_analytical_mf.py::test_pure_kitaev_hamiltonian_and_eigenvalues)

We formulate the microscopic Hamiltonian exclusively in terms of dimensionless Pauli matrices $\sigma_i^\gamma$ ($\gamma \in \{x, y, z\}$):
\begin{equation}
\label{eq:phase4_canonical_hamiltonian}
\mathcal{H} = - K \sum_{\langle i, j \rangle_\gamma} \sigma_i^\gamma \sigma_j^\gamma - \sum_{i} \bm{h} \cdot \bm{\sigma}_i, \qquad K \equiv 1.0, \quad \bm{h} = \frac{h}{\sqrt{3}}\begin{pmatrix} 1 \\ 1 \\ 1 \end{pmatrix},
\end{equation}
such that the Euclidean magnitude of the magnetic field vector is precisely $\|\bm{h}\| = h$.

\begin{table}[htbp]
\centering
\caption{Universal dictionary translating alternative literature conventions to the Phase 4 canonical Pauli standard ($K \equiv 1.0, \|\bm{\sigma}_i\| = 1$).}
\label{tab:phase4_dictionary}
\vspace{2mm}
\begin{tabular}{lccc}
\toprule
\textbf{Physical Observable / Parameter} & \textbf{Phase 4 Canonical} & \textbf{Spin-1/2 Convention} & \textbf{Scaling Ratio} \\
\midrule
Exchange Coupling & $K_{\mathrm{Pauli}} \equiv 1.0$ & $K_{\mathrm{spin}} \equiv 4.0$ & $K_{\mathrm{spin}} = 4 K_{\mathrm{Pauli}}$ \\
Zeeman Field Vector & $\bm{h}_{\mathrm{Pauli}}$ & $\bm{h}_{\mathrm{spin}} \equiv 2 \bm{h}_{\mathrm{Pauli}}$ & $\bm{h}_{\mathrm{spin}} = 2 \bm{h}_{\mathrm{Pauli}}$ \\
Dimensionless Field Ratio & $h / K$ & $(h/K)_{\mathrm{spin}} = \frac{1}{2} (h/K)_{\mathrm{Pauli}}$ & Ratio halved \\
Ground Energy per Site ($h=0$) & $-0.78729862 |K|$ & $-0.19682465 K_{\mathrm{spin}}$ & $E_{\mathrm{spin}} = \frac{1}{4} E_{\mathrm{Pauli}}$ \\
Single-Particle Majorana Energy & $\varepsilon_n = 4 \mathcal{E}_n$ & $\varepsilon_n^{\mathrm{spin}} = \mathcal{E}_n$ & $\varepsilon_{\mathrm{phys}} = 4 \mathcal{E}_{\mathrm{matrix}}$ \\
Low-Field Dirac Gap Asymptotics & $\Delta_{\mathrm{canonical}} = 10.6536 h^3$ & $\Delta_{\mathrm{spin}} = 2.6634 h^3$ & Factor of 4 \\
Pure-Kitaev Matter Hopping & $u_0 = -\bar{U} = -0.5248657$ & $\chi_0 = -\frac{1}{4} u_0 = 0.1312164$ & $u_0 = -4 \chi_0$ \\
\bottomrule
\end{tabular}
\end{table}

\begin{remark}[Code Sign Invariant]
In the project codebase, the exchange bilinear is defined as $\mathcal{H} = \sum_{\langle ij \rangle_\gamma} J \sigma_i^\gamma \sigma_j^\gamma$. Therefore, the ferromagnetic model ($K = +1.0$) is parameterized by $K_{\mathrm{param}} = J = -K = -1.0$. Consequently, the ground-state matter hopping expectation is negative:
\begin{equation}
u_0 = -\bar{U} = -0.5248657053.
\end{equation}
Initializing with $u_0 > 0$ under $K_{\mathrm{param}} = -1.0$ converges to an unphysical inverted-band saddle point where the Dirac gap is an order of magnitude smaller.
\end{remark}

% EDITORIAL QUESTION: How is scientific evidence classified to avoid over-interpreting unprojected mean-field artifacts?
\section{Epistemic Taxonomy and Intermediate Denominators}
\label{sec:phase4_epistemic}

To maintain rigorous transparency, every assertion in Phase 4 is stamped with one of four epistemic tiers:
\begin{enumerate}
    \item \textbf{\texttt{[Proven Theorem]}:} Analytical derivations established by exact operator algebra or group representation theory.
    \item \textbf{\texttt{[Hypothesis Contract]}:} Physical scaling conjectures with pre-registered quantitative verification and refutation criteria.
    \item \textbf{\texttt{[Numerical Benchmark]}:} High-precision floating-point results certified across independent algorithms (e.g., DIIS, Newton-Krylov, and Lanczos).
    \item \textbf{\texttt{[Literature Conjecture]}:} External claims from publications that require independent scrutiny and unit harmonization.
\end{enumerate}

\subsection{Physical Intermediate Resolvents vs. Bare Matrix Hoppings}
In perturbation theory, the effective low-energy Hamiltonian is generated by virtual transitions to high-energy intermediate states via the resolvent $\hat{R} = (E_0 - \hat{H}_0)^{-1}$. 
\begin{itemize}
    \item \textbf{In the Exact Spin Space (Kitaev 2006):} The intermediate state created by a third-order spin operator $\sigma_i^x \sigma_j^y \sigma_k^z$ is an isolated pair of localized $\mathbb{Z}_2$ vortices (visons). In Kitaev's static approximation, the energy denominator is approximated by the bare two-vison threshold:
    \begin{equation}
    \Delta E_v^{(0)} = 2 \Delta_{\mathrm{flux}} \approx 0.2641 |K| \implies (\Delta E_v^{(0)})^2 \approx 0.06975 |K|^2.
    \end{equation}
    \item \textbf{In the Mean-Field Majorana Space:} The intermediate state is a gauge Majorana excitation propagating through the unprojected gauge band. The bare energy denominator is the gauge hopping gap:
    \begin{equation}
    \Delta E_{\mathrm{gauge}}^{(0)} = 2 \bar{U} |K| \approx 1.0497 |K| \implies (\Delta E_{\mathrm{gauge}}^{(0)})^2 \approx 1.1019 |K|^2.
    \end{equation}
\end{itemize}
The ratio of these unrelaxed bare denominators is $(1.0497 / 0.2641)^2 \approx 15.81$. As we prove constructively in Chapter~\ref{ch:phase4_analytical_mf}, this apparent discrepancy is resolved by self-consistent vertex dressing and the dispersive screening of the matter continuum.
